\section{Discussion}
\label{sec:discuss}

\subsection{When Does LLM Prediction Beat Markets?}

Our experiments suggest a nuanced answer to this question that resists both uncritical enthusiasm and blanket skepticism. LLM-augmented prediction provides genuine edge under a specific combination of conditions: the event domain must contain structured information that is publicly available but not yet incorporated into market prices, the LLM must be equipped with current data through context injection, and the market must be sufficiently thin that informed capital has not already eliminated the mispricing.

These conditions are met most consistently in awards prediction (where precursor ceremony results provide strong signals that casual traders may not systematically analyze), commodity pricing (where real-time price gaps between the LLM's stale knowledge and current markets create systematic bias that context injection corrects), and low-attention geopolitical events (where expert knowledge of escalation patterns is not widely distributed among retail traders).

Conversely, high-liquidity sports markets exhibit none of these conditions. Professional bookmakers and quantitative sports bettors have decades of infrastructure for incorporating injury reports, team form, schedule density, and matchup history into odds. Our system provides no informational advantage in these markets, and our experiments confirm negligible improvement from context injection in sports domains.

\subsection{The Alpha Decay Problem}

A fundamental challenge for any prediction system is the dynamic nature of market efficiency. Unlike supervised learning problems where the data distribution is stationary, prediction markets adapt to participant strategies. If Murmur's predictions become publicly available and attract capital, the very mispricings it exploits would disappear as other traders incorporate the same information.

We argue that this concern, while theoretically valid, is practically manageable for three reasons. First, at the capital scales relevant to individual users (\$200--\$10{,}000), order sizes are negligible relative to market depth and do not measurably move prices. Second, the alpha source is not a fixed signal but a \emph{process}---the Context Planner continuously discovers new domain-specific context relevant to emerging events, and the Evolution Engine continuously adapts strategies to the changing efficiency landscape. Third, prediction markets continually generate new events with fresh information asymmetries, providing a renewable source of potential edge even as individual mispricings are corrected.

Nevertheless, we emphasize that our current results demonstrate only marginal improvement over market consensus (Brier Score 0.0934 vs.\ 0.0936), and the statistical significance of this improvement requires larger sample sizes than our current corpus provides. We caution against interpreting these results as evidence of reliable profit-generating alpha and instead frame them as evidence that structured context injection and evolutionary strategy discovery represent promising research directions for prediction market analysis.

\subsection{Architectural Generality}

A distinctive feature of Murmur's design is its architectural generality. The three core mechanisms---context-aware agent reasoning, evolutionary strategy discovery, and agent-based simulation---are not prediction-market-specific. In companion work, the same engine architecture powers LevyFly, a supply chain simulation system where agents make inventory decisions rather than probability estimates, and where the Evolution Engine optimizes reorder policies rather than prediction strategies. The Context Planner in LevyFly injects demand forecasts and supply disruption signals rather than news and prices, but the pipeline structure---domain classification, gap identification, targeted retrieval, structured injection---remains identical.

This architectural convergence suggests that context-enriched agent-based simulation may constitute a general framework for decision-making under uncertainty, with domain-specific instantiation requiring only the definition of appropriate context types, agent behaviors, and fitness functions. We leave formal exploration of this hypothesis to future work.

\subsection{Limitations}

Several limitations qualify our findings. First, the evolved strategy's Brier Score improvement (0.0934 vs.\ 0.0936) is within the margin of statistical noise for our current sample size; significantly more resolved events are needed to establish reliable outperformance. Second, our context injection relies on publicly available information (news, prices, precursor results); access to proprietary data sources---such as satellite imagery of oil storage facilities or internal polling data---could provide substantially larger informational advantages but is beyond our current scope. Third, the Virtual Prediction Market has not yet been fully evaluated, and the agent behavior distributions used in simulation are calibrated from limited historical data. Fourth, our evaluation is specific to Polymarket and may not transfer to other prediction platforms with different market microstructures, fee structures, or participant demographics.

Finally, we note an inherent tension between academic evaluation and practical deployment. Publishing prediction strategies and their performance characteristics may itself erode any edge those strategies provide, as other market participants can read and replicate our approach. This tension is not unique to our work but reflects a broader challenge in financial AI research.
